\documentclass[sigconf]{acmart}
\usepackage{graphicx}
\usepackage{float}
\usepackage{placeins}
\usepackage{amsmath}
\usepackage{graphicx}

\graphicspath{{./}}

\begin{document}
\title{Generative Adversarial Networks in Low Data Regimes}
\author{I. Goodfellow}
\author{Team}
\begin{abstract}
This note summarizes training behavior of adversarial objectives.
\end{abstract}

\setcopyright{none}
\settopmatter{printacmref=false, printccs=false}
\renewcommand\footnotetextcopyrightpermission[1]{}
\maketitle
\thispagestyle{plain}
\pagestyle{plain}
\section{Introduction}
GANs train generator and discriminator simultaneously.

\FloatBarrier
\section{Objective}
The minimax game is
\begin{equation}
\min_G \max_D \mathbb{E}_{x\sim p_{data}}[\log D(x)] + \mathbb{E}_{z\sim p(z)}[\log(1-D(G(z)))].
\end{equation}

\FloatBarrier
\section{Results}
Small-scale tests recover plausible textures and edges.

\FloatBarrier
\section{Conclusion}
Stability depends on architecture and optimization details.

\FloatBarrier
\bibliographystyle{ACM-Reference-Format}
\bibliography{references}
\end{document}
